%%%%%%%%%%%%
%
% $Autor: Wings $
% $Datum: 2019-03-05 08:03:15Z $
% $Pfad: Specifications.tex $
% $Version: 4250 $
% !TeX spellcheck = en_GB/de_DE
% !TeX encoding = utf8
% !TeX root = manual 
% !TeX TXS-program:bibliography = txs:///biber
%
%%%%%%%%%%%%

\chapter{Spezifikationen}

\section{Hardware-Spezifikation}

\renewcommand{\arraystretch}{1.4}

\begin{table}[ht]
	\centering
	\begin{tabular}{|p{6cm}|p{8cm}|}
		\hline
		\textbf{Spezifikation} & \textbf{Beschreibung} \\ \hline
		
		Versorgungsspannung & 230 V AC \\ \hline
		
		Netzanschluss & Kaltgerätekabel (IEC C13) \\ \hline
		
		Schnittverfahren & Heißdraht-Schneiden \\ \hline
		
		Geeignete Materialien & Styropor \\ \hline
		
		Arbeitsbereich &  x mm $\times$ x mm \\ \hline
		
		Abmessungen & L $\times$ B $\times$ H \\ \hline
		
		Gewicht & x kg \\ \hline
		
		Bedienelemente & Start-, Pause- und Stopp-Knopf \\ \hline
		
		Anzeigeeinheit & Integriertes Display zur Status- und Fortschrittsanzeige \\ \hline
		
	\end{tabular}
\end{table}

\section{Software-Spezifikation}

\renewcommand{\arraystretch}{1.4}

\begin{table}[ht]
	\centering
	\begin{tabular}{|p{6cm}|p{8cm}|}
		\hline
		\textbf{Spezifikation} & \textbf{Beschreibung} \\ \hline
		
		Steuerungssystem & CNC-Steuerung über Mikrocontroller \\ \hline
		
		Kommunikation & Verbindung zum Computer über USB \\ \hline
		
		Steuersoftware & [Software] zur Übertragung von Schneiddateien \\ \hline
		
		Dateiformat & G-Code zur Beschreibung der Schneidpfade \\ \hline
		
		Funktionsumfang & Verarbeitung und Ausführung von Schneidbefehlen \\ \hline
		
	\end{tabular}
\end{table}